\label{ch:sterile-implementation}

This chapter maps every equation of \cref{ch:sterile-theory} onto the
concrete \texttt{tpeanuts} implementation, module by module. As with NSI
(see the companion \texttt{BSM\_NSI} document), there is no separate
$N=4$-specific Hamiltonian-assembly module: \pyfunc{hamiltonian\_reduced}
and \pyfunc{hamiltonian\_flavour} in \texttt{core.common.hamiltonian}
handle the 3+1 extension by reading the flavour count directly off
\pyclass{OscillationParameters.pmns}, with no explicit branch on ``is this
sterile.''

\section{\texttt{PMNSSterileParams} and \texttt{PMNS\_sterile}}

\modulehead{core/BSM/bsm\_sterile.py}{\pyclass{PMNS\_sterile}: the
$4\times4$ mixing matrix for one additional light sterile neutrino.}

\begin{implbox}
\pyclass{PMNSSterileParams} is an immutable dataclass holding \emph{only}
the sterile-sector extension: $\theta_{14},\theta_{24},\theta_{34}$ and
their CP phases $\delta_{14},\delta_{24},\delta_{34}$ (the last always
$0$, \cref{ch:sterile-theory}). The Standard Model sector
($\theta_{12},\theta_{13},\theta_{23},\delta_{13}$) lives in a separate,
already-existing \pyclass{PMNSParams} object, passed alongside it to the
\pyclass{PMNS\_sterile} constructor -- mirroring exactly how
\pyclass{NSIConfig} is an optional, orthogonal attribute of
\pyclass{OscillationParameters} rather than baked into the mixing object.
$\Delta m^2_{41}$ is deliberately \emph{not} stored on
\pyclass{PMNSSterileParams}: like $\Delta m^2_{21}$/$\Delta m^2_{3\ell}$,
it never enters a mixing-matrix rotation, so it lives on
\pyclass{OscillationParameters} (via \pyclass{MassSpectrum\_BSM}, below)
instead.

\pyclass{PMNS\_sterile} is a direct subclass of the abstract
\pyclass{core.common.pmns.PMNS} with the class attribute
\code{n\_flavours=4}: every method that is generic in flavour count
(\pyfunc{R12}, \pyfunc{R13}, \pyfunc{R23}, \pyfunc{Delta},
\pyfunc{flavour\_basis}, \pyfunc{select\_antinu}, \dots) is inherited
\emph{unchanged}, automatically sized to $4\times4$; only the genuinely
$3{+}1$-specific pieces are defined here.
\end{implbox}

\begin{theorybox}
The three new rotation builders each act on a single active--sterile
2-plane, mirroring \pyfunc{R12}/\pyfunc{R13}/\pyfunc{R23}'s construction
exactly (via the shared \pyfunc{PMNS.\_rot} generator):
\begin{align*}
  \pyfunc{R14()} &= \text{rotation in the }(e,s)=(0,3)\text{ plane, carrying }\delta_{14}, \\
  \pyfunc{R24()} &= \text{rotation in the }(\mu,s)=(1,3)\text{ plane, carrying }\delta_{24}, \\
  \pyfunc{R34()} &= \text{rotation in the }(\tau,s)=(2,3)\text{ plane, real (}\delta_{34}=0\text{)}.
\end{align*}
\pyfunc{outer\_block} and \pyfunc{reduced} are overridden with the $3{+}1$
product structure of \cref{eq:u4}, identifying $O_4\equiv R_{23}\Delta_4R_{24}R_{34}$
(\pyfunc{PMNS\_sterile.outer\_block}) so that a formal reduced factor
$U_{\rm red,4}\equiv R_{13}\Delta_4^\dagger R_{12}R_{14}$ satisfies
$U_4=O_4\,U_{\rm red,4}$ exactly, mirroring the 3-flavour Standard Model
factorisation. Exactly as for $\Delta$ in the 3-flavour case, $\Delta_4$
cancels from the reduced \emph{kinetic term} specifically,
\begin{equation}
  \keyresult{
    \Delta_4^\dagger\bigl(R_{12}R_{14}\diag(k)R_{14}^\dagger R_{12}^\dagger\bigr)\Delta_4
    = R_{12}R_{14}\diag(k)R_{14}^\dagger R_{12}^\dagger
  },
\end{equation}
valid because neither $R_{12}$ (the $e$--$\mu$ plane) nor $R_{14}$ (the
$e$--$s$ plane) mixes with the $\tau$ index, the only non-trivial index of
$\Delta_4$. This is why \pyfunc{PMNS\_sterile.reduced} returns the
simplified, $\Delta_4$-free
\begin{equation}
  U_{\rm red,4}^{\rm code}=R_{13}R_{12}R_{14}
\end{equation}
rather than the formal $U_{\rm red,4}$: the two agree once contracted with
the diagonal kinetic term, but $U_{\rm red,4}^{\rm code}$ is cheaper to
build and use.
\end{theorybox}

\begin{notebox}
\textbf{$U=O\,U_{\rm red}$ holds exactly only for the formal
$U_{\rm red}$, not for \pyfunc{reduced}.} Consequently,
\pyfunc{outer\_block}$(\,)\cdot$\pyfunc{reduced}$(\,)$ is \emph{not}
bit-for-bit equal to \pyfunc{pmns\_matrix}$(\,)$: the two agree up to a
diagonal phase ($\Delta_4$) that leaves the kinetic term invariant but is
otherwise dropped. This is harmless for every use of \pyfunc{reduced} in
this project (always as the sandwiching matrix of a kinetic term), and
mirrors the identical, already-documented situation for the plain
3-flavour \pyclass{PMNS\_SM}.
\end{notebox}

\section{\texttt{MassSpectrum\_BSM}: the sterile mass splitting}

\modulehead{core/BSM/bsm\_mass\_spectrum.py}{Extends the active-sector
mass-squared-difference vector with $\Delta m^2_{41}$.}

\begin{implbox}
\pyclass{MassSpectrum\_BSM} is a frozen dataclass carrying one extra field,
\code{DeltamSq41}, on top of the inherited \code{DeltamSq21}/\code{DeltamSq3l}.
Its \pyfunc{difference\_vector()} override appends \code{DeltamSq41} to the
base 3-vector, returning a $(\dots,4)$-shaped tensor consumed by
\pyfunc{kinetic\_eigenvalue\_vector} (\cref{sec:impl-kinetic}) to build the
four dimensionless kinetic eigenvalues $k_1,\dots,k_4$. \code{DeltamSq41=None}
is accepted at construction time (so an incomplete configuration can still
be built and inspected) but raises \code{ValueError} the moment
\pyfunc{difference\_vector} is actually evaluated, since a 3+1 Hamiltonian
genuinely cannot be assembled without it.
\end{implbox}

\section{Hamiltonian assembly for $N=4$}
\label{sec:impl-kinetic}

\modulehead{core/common/hamiltonian.py}{\pyfunc{hamiltonian\_kinetic\_reduced}/
\pyfunc{hamiltonian\_matter\_reduced}: the same two functions used for the
3-flavour and NSI cases, generalised to any flavour count.}

\begin{theorybox}
The reduced-basis kinetic term uses $U_{\rm red,4}^{\rm code}=R_{13}R_{12}R_{14}$
together with its \emph{Hermitian conjugate} (never the plain transpose):
\begin{equation}
  \keyresult{H_{\rm kin,4}^{\rm red} = U_{\rm red,4}^{\rm code}\,\diag(k_1,k_2,k_3,k_4)\,\bigl(U_{\rm red,4}^{\rm code}\bigr)^\dagger}.
\end{equation}
Unlike the 3-flavour $R_{13}R_{12}$, $U_{\rm red,4}^{\rm code}$ is
genuinely complex whenever the active--sterile CP phase $\delta_{14}\neq0$
(it carries $R_{14}$, built with a non-trivial phase, unlike
$R_{12}$/$R_{13}$): the Hermitian-conjugate convention is therefore not a
merely-safe default here but the only correct choice for
$U_4=O_4\,U_{\rm red,4}$ to be respected by the flavour-basis dressing.
\end{theorybox}

\begin{theorybox}
The matter term takes one of three forms, matching
\cref{sec:sterile-matter}'s three cases exactly:
\begin{enumerate}
  \item \textbf{No NSI, no neutral-current term (the default).} Since
    $R_{23},\Delta_4,R_{24},R_{34}$ (i.e.\ every generator of $O_4$) leave
    the electron index fixed, $O_4\,\diag(V_{CC},0,0,0)\,O_4^\dagger=\diag(V_{CC},0,0,0)$
    exactly, so no rotation is needed at all:
    $H_{\rm mat,4}^{\rm red}=H_{\rm mat,4}^{\rm flavour}=\diag(V_{CC},0,0,0)$.
    This introduces no approximation beyond the 3-flavour case, for any
    $\theta_{14},\theta_{24},\theta_{34}$.
  \item \textbf{With the neutral-current term, no NSI.} Supplying a
    neutron density activates $H_{\rm mat,4}^{\rm flavour}=\diag(V_{CC},0,0,-V_{NC})$
    (\cref{eq:sterile-nc}) instead. Because $O_4$ genuinely mixes the
    $\{\mu,\tau,s\}$ block, it no longer leaves this term invariant, so
    \pyfunc{hamiltonian\_matter\_reduced} takes the general rotated path,
    \begin{equation}
      \keyresult{H_{\rm mat,4}^{\rm red}=O_4^\dagger\,H_{\rm mat,4}^{\rm flavour}\,O_4},
    \end{equation}
    even though NSI is off -- the reduction-trick shortcut of case~1 is
    specific to the charged-current-only term.
  \item \textbf{With NSI (and, optionally, the neutral-current term).}
    Fully supported by the same code path documented in the companion
    \texttt{BSM\_NSI} document: \pyfunc{NSIConfig.epsilon\_tensor(n\_flavours=4)}
    (and \pyfunc{epsilon\_n\_tensor}) embed the $3\times3$ active-flavour
    block in the top-left corner of a $4\times4$ zero matrix, since $\nu_s$
    is a gauge singlet and carries no NSI coupling of either kind; the
    resulting flavour-basis matter term is rotated into the reduced basis
    with the same $O_4^\dagger(\cdot)O_4$ transformation as case~2.
\end{enumerate}
Nothing in \pyfunc{hamiltonian\_reduced}/\pyfunc{hamiltonian\_flavour}
branches explicitly on ``sterile or NSI or both'': the flavour count
$n\in\{3,4\}$ and the presence of \code{oscillation.nsi} are read directly
off \pyclass{OscillationParameters}, so all combinations are handled by the
exact same two functions used throughout this project.
\end{theorybox}

\section{The closed-form quartic (Ferrari) eigenvalue solver}
\label{sec:impl-ferrari}

\modulehead{core/perturbative/spectral.py}{\pyfunc{hamiltonian\_traceless\_eigenvalues(\dots,
analytic=True)}: closed-form eigenvalues for a traceless Hermitian
Hamiltonian, Cardano for $N=3$, Ferrari for $N=4$.}

\begin{theorybox}
The perturbative evolutor (\cref{sec:impl-perturbative-sterile}) needs the
Hamiltonian's spectral decomposition $H=\sum_a\lambda_a M_a$. By default
$\lambda_a$ is obtained numerically via \pyfunc{torch.linalg.eigvalsh}; an
explicit closed-form alternative is also available specifically for
$N\in\{3,4\}$, since a traceless Hermitian matrix's characteristic
polynomial is guaranteed real-rooted and has no sub-leading term. For
$N=4$ specifically -- the case this document is about -- the characteristic
polynomial of the traceless part $T=H-\Tr(H)I_4/4$ is the depressed
quartic $\lambda^4+p\lambda^2+q\lambda+r=0$ with $p=c_1=-\tfrac12\Tr(T^2)$,
$q=-e_3=-\tfrac13\Tr(T^3)$, $r=e_4=\det(T)$. Ferrari's method solves it by
first finding a real root $m$ of the \emph{resolvent cubic}
\begin{equation}
  m^3+p\,m^2+\Bigl(\tfrac{p^2}4-r\Bigr)m-\tfrac{q^2}8=0,
\end{equation}
itself solved with the same trigonometric (Cardano) formula used for the
$N=3$ case: since $T$ is Hermitian, the quartic is guaranteed real-rooted,
and the resolvent cubic's three roots are exactly the real pairwise sums
$m_{ij|kl}=-\tfrac12(\lambda_i+\lambda_j)(\lambda_k+\lambda_l)$ over the
three ways of splitting the four roots into two pairs, so it is real-rooted
too and the closed form applies unconditionally. Taking the \emph{largest}
resolvent root $m$ (the conditioning-optimal choice) factors the quartic
into two real quadratics,
\begin{equation}
  \keyresult{
    s=\sqrt{2m},\qquad
    \lambda\in\Bigl\{\tfrac{s\pm\sqrt{\Delta_1}}2,\ \tfrac{-s\pm\sqrt{\Delta_2}}2\Bigr\},
    \qquad
    \Delta_1=-2\Bigl(m+p+\tfrac qs\Bigr),\ \ \Delta_2=-2\Bigl(m+p-\tfrac qs\Bigr).
  }
\end{equation}
\end{theorybox}

\begin{notebox}
\textbf{Conditioning is not unconditional for $N=4$, unlike $N=3$.} The
chosen pairing's $m$ can itself be close to singular (only possible, for a
traceless Hermitian $T$, near $T\approx0$), or the resolvent cubic's own
three roots can nearly coincide even when the \emph{chosen} $m$ is
well-separated from zero (e.g.\ $\lambda=(a,a,-a,-a)$ has a
well-conditioned pairing $m=2a^2$ whose resolvent cubic nonetheless has a
repeated root elsewhere). Batch entries flagged by either check fall back
directly to \pyfunc{torch.linalg.eigvalsh} for the \emph{eigenvalues}
themselves (unlike the projector-only degeneracy fallback described next),
since this closed form can lose the eigenvalues, not merely lose precision
in a downstream ratio.

\textbf{Spectral projectors and the degeneracy fallback.} The corresponding
$N=4$ projector formula,
\begin{equation}
  M_a = \frac{T^3+\lambda_aT^2+(c_1+\lambda_a^2)T+(\lambda_a^3+c_1\lambda_a-e_3)I_4}{4\lambda_a^3+2c_1\lambda_a-e_3},
\end{equation}
has a denominator equal to $p'(\lambda_a)$, the characteristic polynomial's
derivative at that root, which vanishes exactly at a repeated eigenvalue.
Near a (near-)degenerate spectrum, \pyfunc{\_spectral\_degeneracy\_mask}
flags the affected batch entries (by relative eigenvalue-gap size) and
replaces \emph{only their projectors} wholesale with
$M_a=v_av_a^\dagger$ from a direct \pyfunc{torch.linalg.eigh(T)} call --
exact and stable by construction, computed only when at least one batch
entry actually needs it.
\end{notebox}

\begin{implbox}
\pyfunc{hamiltonian\_traceless\_eigenvalues(T, *, analytic=False, c1=None,
c0=None, e3=None, e4=None)} dispatches on \code{T.shape[-1]}: Cardano for
$3$, Ferrari for $4$, and unconditional fallback to
\pyfunc{torch.linalg.eigvalsh} for any other size (no closed form is
defined there). \pyfunc{evolutor\_perturbative\_segment} and
\pyfunc{earth\_evolutor}/\pyfunc{atmosphere\_evolutor} expose this as the
public \code{analytic\_eigenvalues=True} flag; passing it together with a
matrix-exponential (\code{method="numerical"}) propagation path raises,
since that path never diagonalises the Hamiltonian at all.
\end{implbox}

\section{The perturbative segment evolutor}
\label{sec:impl-perturbative-sterile}

\modulehead{core/perturbative/evolutor.py}{\pyfunc{evolutor\_perturbative\_segment}:
generalised, without modification, to $N=4$.}

\begin{implbox}
\pyfunc{evolutor\_perturbative\_segment} never branches explicitly on
``sterile'': it reads \code{n\_flavours = int(pmns.n\_flavours)} off
whatever \pyclass{PMNS} object \pyclass{OscillationParameters} carries and
sizes every subsequent tensor accordingly. The one behaviour that
\emph{is} sterile-specific is \code{include\_matter\_nc}: when True (and
\code{n\_flavours==4}), the sterile neutral-current direction matrix
\code{P\_nc} (\cref{eq:sterile-nc}, rotated by $O_4^\dagger(\cdot)O_4$
exactly as in \cref{sec:impl-kinetic}'s case~2) is built by
\pyfunc{\_matter\_direction\_reduced} and sandwiched into the first-order
correction alongside the CC direction \code{P}, using the profile's
\code{residual\_integral\_neutron(la, lb)} and
\pyfunc{matter\_potential\_nc}. Requesting it without a profile built with
neutron-density coefficients (\code{potential\_n}/\code{average\_n}/
\code{residual\_integral\_neutron} unavailable) raises \code{ValueError}
rather than silently omitting the term.
\end{implbox}

\section{Pipeline wiring: the neutral-current policy}

\begin{notebox}
\textbf{\code{resolve\_include\_matter\_nc}: the shared tri-state policy.}
Because omitting the sterile NC term is a genuine physical incompleteness
once $\theta_{24}$ or $\theta_{34}$ are non-zero (\cref{sec:sterile-matter}),
\texttt{medium.solar}, \texttt{medium.earth}, and \texttt{medium.atmosphere}
all resolve \code{include\_matter\_nc} through one shared function in
\texttt{core.common.oscillation}:
\begin{itemize}
  \item \code{include\_matter\_nc} given explicitly (\code{True}/\code{False}):
    always honoured as-is.
  \item \code{include\_matter\_nc=None} (unspecified): resolves to
    \code{True} when \code{oscillation.BSM\_extension\_sterile} is set
    \emph{and} the profile/medium in hand can actually supply a neutron
    density; \code{False} for the plain 3-flavour case regardless (there
    $V_{NC}$ truly is an unobservable common phase, so omitting it is
    exact, not approximate). When sterile is active but no neutron-density
    data is available, this still resolves to \code{False} (the CC-only
    fallback), but emits an explicit \code{RuntimeWarning} rather than
    resolving silently -- exactly the warning quoted in
    \cref{sec:results-intrinsic3}, triggered because the
    \texttt{even\_power} Earth profile used by that notebook's Section~5
    was not built with \code{include\_neutron=True}.
\end{itemize}
Each medium then supplies the raw density itself: \pyclass{EarthProfile}
via \code{profile\_perturbative\_kwargs={'include\_neutron': True}} at
construction time (PREM- or even-power-tabulated), \pyclass{AtmosphereParameters}
via its own \code{include\_matter\_nc} field feeding
\pyfunc{atmosphere\_density(\dots, density\_type="neutron\_density")}, and
\pyclass{SolarMediumProfile} via its \code{density\_n} table (derived from
the standard solar model composition).
\end{notebox}

\section{Summary: modules touched by the sterile extension}

\begin{center}
\begin{tabular}{@{}ll@{}}
\toprule
Module & Role \\
\midrule
\pkgpath{core/BSM/bsm\_sterile.py} & \pyclass{PMNS\_sterile}: $4\times4$ mixing matrix, $R_{14}/R_{24}/R_{34}$ \\
\pkgpath{core/BSM/bsm\_mass\_spectrum.py} & \pyclass{MassSpectrum\_BSM}: appends $\Delta m^2_{41}$ \\
\pkgpath{core/common/hamiltonian.py} & \pyfunc{hamiltonian\_reduced}/\code{\_flavour}, generic in flavour count \\
\pkgpath{core/perturbative/spectral.py} & Ferrari closed-form quartic eigenvalues, $N=4$ \\
\pkgpath{core/perturbative/evolutor.py} & \code{P\_nc} sandwich term, \code{analytic\_eigenvalues} \\
\pkgpath{core/numerical/evolutor.py} & Exact matrix-exponential evolution, generic in flavour count \\
\pkgpath{core/common/oscillation.py} & \pyfunc{resolve\_include\_matter\_nc}: shared tri-state policy \\
\pkgpath{medium/earth}, \pkgpath{medium/atmosphere}, \pkgpath{medium/solar} & Neutron-density profile construction/sampling \\
\bottomrule
\end{tabular}
\end{center}
